\chapter{Stabilize or explore: A circuit switch for odor plume tracking}

\section{Introduction}

In the natural world, odor often disperses into elongated plumes composed of packets of fine filaments \citep{murlis2000, murlis1992}. % TODO: murlis2000 = Murlis et al. 2000; murlis1992 = Murlis et al. 1992
When the wind is blowing steadily in a relatively constant direction, and animal can navigate toward an odor source by moving in the upwind direction \citep{vanbreugel2014, willis2005}; % TODO: vanbreugel2014 = van Breugel and Dickinson 2014; willis2005 = Willis and Avondet 2005
visual cues can also help orient an animal toward an odor source \citep{vanbreugel2014, duistermars2008}. % TODO: duistermars2008 = Duistermars and Frye 2008
But even in the absence of wind cues or visual cues, an animal can use the spatio-temporal statistics of the odor to align its path to an odor plume \citep{vickers1994, krishnan2011} --- a behavior we call ``plume tracking''. % TODO: vickers1994 = Vickers and Baker 1994; krishnan2011 = Krishnan et al. 2011
Plume tracking using odor statistics alone is a difficult inference problem, because odor packets can occur outside the plume, and odor encounters are intermittent even inside the plume. The inside and outside of the plume are potentially discriminable based the frequency of odor encounters, but these encounters are irregular and intermittent in both regimes. Many insects solve the plume-tracking problem using a simple algorithm: they move forward in the direction that maintains a high odor encounter frequency, and they turn when odor encounter frequency drops too low \citep{vickers1994, krishnan2011}. % TODO: vickers1994 = Vickers and Baker 1994; krishnan2011 = Krishnan et al. 2011
This algorithm also emerges spontaneously in artificial ``agents'' that have been asked to follow odor plumes \citep{singh2023}. % TODO: singh2023 = Singh et al. 2023

This algorithm must require a process that integrates odor encounters over time to estimate the current odor encounter frequency. It must also require a process that can adjust path straightness, based on this estimate of odor encounter frequency. In insects, both these processes have been localized to discrete modules in the brain: olfactory learning is localized to the mushroom body, whereas the control of navigation is localized to the central complex. Thus, it is logical to hypothesize that the pathway that connects the mushroom body to the central complex mediates plume tracking based on odor statistics.

In \textit{Drosophila}, recent connectome analyses have revealed that there are a few pathways that connect mushroom body output neurons (MBONs) to the central complex \citep{li2020, hulse2021}. % TODO: li2020 = Li et al. 2020; hulse2021 = Hulse et al. 2021
One of the most prominent of these pathways involves two specific mushroom body output neurons: MBON09, which is GABAergic, and MBON21, which is cholinergic \citep{aso2014}. % TODO: aso2014 = Aso et al. 2014
MBON09 and MBON21 converge onto FB4R \citep{li2020, hulse2021}, % TODO: li2020 = Li et al. 2020; hulse2021 = Hulse et al. 2021
a glutamatergic tangential cell of the fan-shaped body \citep{wolff2025}; % TODO: wolff2025 = Wolff et al. 2025
the top outputs of FB4R cells are the h$\Delta$B neurons of the fan-shaped body. h$\Delta$B neurons are interesting because they encode an internal estimate of the fly's travel vector, i.e.\ its translational velocity in allocentric coordinates \citep{lu2021, lyu2021}. % TODO: lu2021 = Lu et al. 2021; lyu2021 = Lyu et al. 2021
Moreover, h$\Delta$B neurons provide strong disynaptic input to the PFL2/3 neurons that influence steering and speed control in walking flies \citep{westeinde2024, musselspires2024}. % TODO: westeinde2024 = Westeinde et al. 2024; musselspires2024 = Mussells Pires et al. 2024
Modulations in this h$\Delta$B signal have the potential to change the fly's path, because they are positioned to adjust the internal setpoint for central complex steering control. Overall, this anatomy suggests a model where MBONs estimate the current odor encounter frequency, and this estimate is used to modulate path straightness, to promote turning outside the odor plume, but straight walking inside the odor plume.

\section{Results}

To study plume tracking behavior in virtual reality, we created a paradigm where odor encounter frequency depends on the direction of locomotion. Specifically, we placed flies on a spherical treadmill which allows them to walk freely in any direction in a virtual environment. To anchor the fly’s head direction system, we placed a visual cue at a fixed azimuthal position in the virtual environment. After we allowed the fly to acclimate to this environment, we switched on the ``virtual plume'': we defined heading range where an odor delivery valve was set to the open state, and we allowed the fly to orient freely relative to this plume. The valve delivered odor into a continuous air stream, but the fly’s antennae were stabilized so that the fly could not detect the direction of airflow; thus, any plume-tracking behavior in this paradigm must rely on odor, not directional wind cues. (Orienting in the upwind direction can be a useful strategy for finding an odor source, but wind cues are not always available; for this reason, we focus in this study on the problem of plume tracking without the benefit of directional wind cues.) 

In this environment, after we switched on the virtual plume, flies tended to reorient their heading toward an odor edge (Fig.~\ref{fig:TODO}). % TODO: add figure label
As their heading moved back and forth across this edge, they experienced a high frequency of their odor encounters. This is clearly an active behavior, because their odor encounter frequency was significantly higher than it would have been, had they not reoriented toward the odor edge (Fig.~\ref{fig:TODO}). % TODO: add figure label
Together, these results indicate that head-fixed flies in a virtual reality environment can actively reorient in a direction that produces a high rate of odor encounters. We term this ``virtual plume tracking''.

Natural odor plumes can be described as having three levels of macroscopic structure: an odor filament (the smallest macroscopic level), a burst (a set of nearby filaments), and a plume (a set of nearby bursts; note that isolated bursts also occur outside the plume). Similarly, we can describe three levels of structure in our virtual environment: an odor encounter (i.e., an odor valve opening - this is the elemental unit of odor delivery, analogous to a filament), a burst (a set of nearby odor encounters), and a virtual plume (a set of nearby bursts). 

In our virtual reality environment, we found that individual odor encounters did not reliably modulate behavior, but odor bursts did. Specifically, odor burst onset increased rotational speed and decreased path straightness. Odor burst onset also decreased forward velocity. Notably, these behavioral responses were significantly weaker inside the virtual plume. As a result, flies were able to walk relatively straight and fast when they were inside the plume. These results suggest there are neural mechanisms to learn the statistics of odor encounters, and to link these statistics to locomotor control. 

Motivated by these results, we looked for neural signals that could link odor statistics with locomotion. Previous work has shown that mushroom body output neurons (MBONs) can learn the statistics of odor encounters, and some MBONs send anatomical output to the fan-shaped body, a region of the central complex which monitors and influences the fly’s locomotor path. In particular, the fan-shaped bdy contains a topographic representation of allocentric travel direction in h$\Delta$B neurons,and h$\Delta$B neurons receive prominent anatomical input from two mushroom body output neurons MBONs, MBON09 and MBON21, via an interposed tangential cell (FB4R). h$\Delta$B neurons in turn provide major disynaptic output to the central complex output neurons that influence steering and forward velocity in walking flies (PFL2/3 neurons). This circuit suggests a mechanism by which odor statistics might influence navigation.

To investigate this mechanism, we performed calcium imaging in MBON09 and MBON21 as flies navigated in our virtual plume tracking paradigm. We expressed the fast calcium indicator jGCaMP7f in these MBONs using specific genetic driver lines,  and we monitored their activity using 2-photon imaging. 

Notably, we found that odors had opposing effects on these two cell types. Outside the plume, odor burst onset produced inhibition of MBON21, but excitation of MBON09. Conversely, inside the plume, both these responses reversed: odor burst onset produced excitation of MBON21, but inhibition of MBON09. These opposing odor responses are interesting because MBON21 is cholinergic (excitatory), while MBON09 is likely inhibitory (glutamatergic). Our results imply that, inside the plume, odor burst onset should drive an increase in excitation (via MBON21) accompanied by a decrease in inhibition (via MBON09). Conversely, outside the plume, odor burst onset should drive an decrease in excitation (via MBON21) accompanied by an increase in inhibition (via MBON09). In short, transitions into and out of the plume should produce marked changes in downstream activity in the fan-shaped body, which could help explain why a fly has different locomotor responses to odors inside versus outside the plume. 

One caveat associated with any virtual reality experiment is that odor and behavioral state are intrinsically coupled, because the fly’s behavior is what controls odor valve opening and odor statistics. To uncouple odor from behavior, we delivered odor in an replay paradigm using data from our virtual plume experiments. In this replay paradigm, odor valve openings taken from the virtual plume assay were replayed to different flies, thereby preserving the temporal statistics of odor encounters while breaking the contingency between odor delivery and the animal’s own behavior. Under these replay conditions, MBON09 still exhibited odor responses that depended on odor statistics. Specifically, low-frequency odor pulses drove strong excitation in MBON09, whereas high-frequency odor pulses drove weak excitation, or even inhibition. Meanwhile, under replay conditions, MBON21…. These results argue that MBON odor responses do indeed depend on the frequency of odor encounters, in a manner consisted with our visual plume experiments. 

We next asked whether MBON odor responses might predict changes in the  activity of h$\Delta$B neurons. We expressed a synaptically targeted calcium indicator (syt-jGCaMP7f) in h$\Delta$B cells, along with either MBON09 or MBON21. In these experiments, we presented replayed odor sequences derived from our virtual plume experiments; we also delivered regularly spaced odor pulses, allowing us to examine MBON–h$\Delta$B relationships across a range of odor statistics. Consistent with prior studies, we found that the population activity of h$\Delta$B neurons has a roughly sinusoidal spatial profile across the horizontal axis of the fan-shaped body; the phase of this sinusoid correlates with the fly’s travel direction, while the amplitude correlates with the fly’s travel speed. We also noted that the mean level of activity in h$\Delta$B neurons is generally correlated with bump amplitude, but these parameters can also vary independently to some degree. 

We found that MBON09 activity was positively correlated with h$\Delta$B mean (strongly) and h$\Delta$B amplitude (more weakly). This positive correlation likely reflects the fact that MBON09 is GABAergic, and FB4R is glutamatergic, and so likely inhibitory; thus, MBON09 activity is likely to disinhibit h$\Delta$B neurons. Because each FB4R neuron synapses onto the entire h$\Delta$B population, MBON09 activity is expected to increase mean h$\Delta$B activity across the fan-shaped body. The h$\Delta$B amplitude might increase as a direct consequence of this global excitation -- e.g., due to any nonlinearily (like the spike threshold) that causes global excitation to amplify the effect of focal excitation.

% TODO: incomplete paragraph — opening bracket below indicates draft text
[Conversely, MBON21 activity was negatively correlated with h$\Delta$B mean and h$\Delta$B amplitude. This likely reflects the fact that MBON21 is cholinergic, and so likely excitatory; thus, its activity is likely to inhibit h$\Delta$B neurons, via recruitment of FB4R.
In these experiments, we also examined the timing of neural dynamics around odor encounters, based on the sign of the MBON odor response. When an odor elicited an increase in MBON09 activity, or a decrease in MBON21 activity, we found that this odor encounter was accompanied by a clear, time-locked increase in the h$\Delta$B mean. This is consistent with the idea that MBON09 excitation and MBON21 inhibition drive depolarization (disinhibition) of the entire h$\Delta$B population. This change in h$\Delta$B activity was not associated with an immediate time-locked change in locomotion, arguing it does not arise from a change in the locomotor-related (non-MBON) inputs to h$\Delta$B. Recall that this MBON response signature (MBON09 increasing, MBON21 decreasing) is typical of odor encounters outside the virtual plume, where odor packets are sparse. In this regime, our results indicate that odors produce a time-locked increase in h$\Delta$B activity, which may contribute to increased exploratory turning to relocate the odor plume.
We also identified odor encounters that % TODO: incomplete sentence
Given that MBON09 is positively correlated with h$\Delta$B mean activity (and bump amplitude), we next asked what MBON/h$\Delta$B activity generally does to goal-directed navigation—i.e., whether these signals act in a way that would stabilize a current travel goal or, instead, promote reorientation away from it. If MBON09 and MBON21 indeed influence goal-directed control through the fan-shaped body, then their activity should be meaningfully related to the animal’s heading representation. Consistent with the prediction, MBON09 and MBON21 both exhibited robust heading tuning, but with opposite phase relationships to the fly’s recent goal: MBON21 activity was highest near the goal direction, whereas MBON09 peaked near the direction opposite the goal. These MBON signals were also modulated by locomotor state, with MBON21 positively correlated with forward velocity and MBON09 positively correlated with yaw speed, but their heading tuning preserves across the locomotion range. 
Motivated by these opponent MBON tunings and their opposite correlations with h$\Delta$B, we then asked how h$\Delta$B activity was organized relative to the internal heading representation. We analyzed mean h$\Delta$B activity and h$\Delta$B bump amplitude as a function of bump phase. Both mh$\Delta$B activity and h$\Delta$B bump amplitude was elevated when the bump phase deviates away from the recent goal phase. In addition, both h$\Delta$B mean activity and bump amplitude increased with forward velocity, while this phase-dependent organization was preserved across velocity conditions. Together, these results suggest that MBON21 activity aligns with—and may help stabilize—goal-directed navigation, while higher MBON09 and correspondingly h$\Delta$B activity are associated with destabilizing and reorienting away from the recent goal direction.
To test whether h$\Delta$B activity plays a causal role in regulating goal-directed navigation, we performed optogenetic activation of h$\Delta$B by expressing the red-light-gated cation channel CsChrimson in these neurons. We then compared heading behavior before, during, and after stimulation in experimental flies and genetic or effector matched controls. We quantified how the fly’s instantaneous heading changed relative to its recent navigational goal, defined as the modal heading over a prolonged history time window.
Activation of h$\Delta$B caused flies to continuously turn away from this recent setpoint direction during the stimulation period. On average, experimental flies showed a clear increase in circular distance from their recent modal heading during stimulation, whereas genetic and effector controls did not show a comparable shift. This effect persisted at least partially into the post-stimulation period, consistent with a transient destabilization of the current heading setpoint rather than a momentary perturbation alone. These results suggest that globally elevating h$\Delta$B activity is sufficient to actively drive flies away from their current navigational setpoint. 
To test the causal contribution of MBON21 to locomotor control, we performed optogenetic activation in head-fixed walking flies using targeted two-photon stimulation. This approach was motivated by the off-target expression present in available MBON21 driver lines and allowed stimulation to be spatially restricted to the recorded target region. Flies co-expressed CsChrimson and jGCaMP8m in MBON21, enabling us to verify in each experiment that stimulation of the experimental site produced a robust calcium response in the target neuron. Within each trial, the laser was directed sequentially to a control brain region, then to the MBON21 target region, and then back to the same control region, with each epoch matched in duration. Behavioral comparisons were therefore made between control- and experimental-stimulation periods within the same animal and time window.
Stimulation of the experimental site reliably evoked a clear increase in MBON21 fluorescence and produced a clear change in locomotor state: stimulation of the target region reduced yaw speed while largely preserving forward velocity on average. Thus, activating MBON21 selectively suppressed turning without causing a general decrease in locomotion. This effect is consistent with the behavioral signature associated with dense odor encounter and edge-tracking states, in which MBON21 is preferentially active and flies exhibit reduced turning together with sustained forward movement. These results therefore support a causal role for MBON21 in stabilizing ongoing heading by biasing locomotion away from reorientation and toward straighter walking.
We next asked whether the odor-history-dependent locomotor signatures measured in real flies are sufficient to produce odor-edge tracking. In the experimental data, odor encounters were associated with structured changes in yaw, lateral, and forward velocity, and these responses depended on recent encounter history, with distinct kinematic adjustments emerging in sparse and dense encounter regimes. To test whether these local response rules could account for the emergent behavior, we incorporated the corresponding history-dependent velocity modulations into the simulation and examined the resulting trajectories after odor onset. Simulated flies then exhibited persistent movement along the odor boundary, closely resembling the odor-edge tracking behavior observed in vivo. To quantify this effect across initial heading conditions and random seeds, we measured both the fraction of time spent in the odor sector and the frequency of sector entries before and after odor onset. Because odor was only available after the imposed onset time, the pre-onset measurements do not reflect true odor sampling, but instead provide a geometric baseline relative to the same sector that subsequently defined odor availability. Relative to this baseline, odor onset increased both sector occupancy and encounter frequency, indicating that activation of odor-history-dependent kinematic responses biased trajectories toward repeated returns to, and prolonged residence near, the odor boundary rather than a transient reorientation alone. Together, these results suggest that odor-history-dependent modulation of locomotor kinematics is sufficient to generate robust odor-edge tracking in the model.

\section{Methods}

\subsection{Flies}

Flies that expressed CsChrimson were raised on cornmeal--molasses food (Archon Scientific) supplemented with 100~\textmu l all-trans-retinal (Sigma; 17~mM in ethanol) and kept wrapped in light-blocking foil to avoid photoconversion. Flies that did not express CsChrimson were raised on standard cornmeal--molasses food. All flies were housed in an incubator under a 12~h:12~h light:dark cycle at 25~\textdegree C and 50--70\% relative humidity.

Experimenters were aware of fly genotype in all cases except for the hDeltaB optogenetic stimulation experiment. \textbf{TODO: add reference to the corresponding main-text figure.} Sample sizes were selected according to field-standard conventions, which are typically based on the anticipated degree of animal-to-animal variability as informed by published studies and pilot experiments. All experiments were performed using flies carrying at least one wild-type copy of the white (\textit{w}) gene.

\paragraph{Genotypes used in each figure.}

{\raggedright
\begin{itemize}
    \item MBON21 calcium imaging: \textit{w}/+; +; \textit{P}\{y[+t7.7] w[+mC]=VT999036-GAL4\}attP2/\textit{PBac}\{y[+t7.7] \textit{PBac}\{y[+mDint2] w[+mC]=20XUAS-IVS-jGCaMP7f\}VK00005.
    \item MBON09 calcium imaging: \textit{w}/+; +; \textit{P}\{y[+t7.7] w[+mC]=R94B10-GAL4.DBD\}attP2 \textit{PBac}\{y[+mDint2] w[+mC]=R20A02-p65.AD\}VK00027/\textit{PBac}\{y[+t7.7] \textit{PBac}\{y[+mDint2] w[+mC]=20XUAS-IVS-jGCaMP7f\}VK00005.
    \item hDeltaB calcium imaging: \textit{w}/+; +/\textit{P}\{R72B05-p65.AD\}; attP40\textit{P}\{y[+t7.7] w[+mC]=UAS-sytjGCaMP7f\}attP2;/\textit{P}\{VT055827-Gal4.DBD\}attP2.
    \item hDeltaB and MBON09 calcium imaging: \textit{w}/+; \textit{P}\{y[+t7.7] w[+mC]=UAS-Syt1::jGCaMP7f\}attP40/\textit{P}\{R72B05-p65.AD\}attP40; \textit{P}\{y[+t7.7] w[+mC]=GMR94B10-GAL4\}attP2/\textit{P}\{VT055827-Gal4.DBD\}attP2.
    \item hDeltaB and MBON21 calcium imaging: \textit{w}/+; \textit{P}\{y[+t7.7] w[+mC]=UAS-Syt1::jGCaMP7f\}attP40/\textit{P}\{R72B05-p65.AD\}attP40; \textit{P}\{y[+t7.7] w[+mC]=VT999036-GAL4\}attP2/\textit{P}\{VT055827-Gal4.DBD\}attP2.
    \item hDeltaB optogenetics: \textit{w}/+; \textit{P}\{y[+t7.7] w[+mC]=20XUAS-IVS-CsChrimson.mVenus\}attP40/\textit{P}\{R72B05-p65.AD\}attP40; \textit{P}\{VT055827-Gal4.DBD\}attP2/+.
    \item MBON21 optogenetics and imaging: \textit{w}/+; \textit{PBac}\{y[+mDint2] w[+mC]=20X-UAS-IVS-CsChrimson::mCyRFP3\}VK00002/+; \textit{P}\{y[+t7.7] w[+mC]=VT999036-GAL4\}attP2/\textit{PBac}\{y[+mDint2] w[+mC]=20XUAS-IVS-Syt1::jGCaMP8m\}VK00005.
    \item MBON09 optogenetics and imaging: \textit{w}/+; \textit{PBac}\{y[+mDint2] w[+mC]=20X-UAS-IVS-CsChrimson::mCyRFP3\}VK00002/+; \textit{P}\{y[+t7.7] w[+mC]=R94B10-GAL4.DBD\}attP2 \textit{PBac}\{y[+mDint2] w[+mC]=R20A02-p65.AD\}VK00027/\textit{PBac}\{y[+mDint2] w[+mC]=20XUAS-IVS-Syt1::jGCaMP8m\}VK00005.
\end{itemize}
}

\paragraph{Origins of transgenic stocks.}

{\raggedright
The following stocks were obtained from the Bloomington Drosophila Stock Center (BDSC) and previously published as follows: \textit{P}\{y[+t7.7] w[+mC]=R94B10-GAL4.DBD\}attP2 \textit{PBac}\{y[+mDint2] w[+mC]=R20A02-p65.AD\}VK00027 (BDSC\_68262), \textit{P}\{y[+t7.7] w[+mC]=20XUAS-IVS-CsChrimson.mVenus\}attP40 (BDSC\_55135), \textit{P}\{y[+t7.7] w[+mC]=UAS-sytjGCaMP7f\}attP2 (BDSC\_94619), \textit{P}\{y[+t7.7] w[+mC]=UAS-sytjGCaMP7f\}attP40 (BDSC\_94620), \textit{PBac}\{y[+mDint2] w[+mC]=20XUAS-IVS-CsChrimson::mCyRFP3\}VK00002 (BDSC\_606323).
The following stock was obtained from Yoshinori Aso at Janelia Research Campus: \textit{P}\{y[+t7.7] w[+mC]=VT999036-GAL4\}attP2.
The following stock was obtained from Gaby Maimon at Rockefeller University: \textit{P}\{R72B05-p65.AD\}attP40; \textit{P}\{VT055827-Gal4.DBD\}attP2.
The following stock was obtained from Thomas R.\ Clandinin at Stanford University: \textit{PBac}\{y[+mDint2] w[+mC]=20XUAS-IVS-Syt1::jGCaMP8m\}VK00005.

}

\subsection{Fly preparation and dissection}

Flies were harvested at least 24 hours before each experiment. For all calcium imaging experiments, they were maintained on standard molasses food, whereas for optogenetic manipulation experiments they were kept on molasses food supplemented with 100~\textmu l all-trans-retinal (Sigma; 17~mM in ethanol). In both types of experiments, only female flies 20--72~h after eclosion were used. For odor-related experiments, flies were food-deprived starting the day before testing and housed with a kimwipe moistened with deionized water to maintain humidity. Experiments were conducted without any circadian timing constraints.

The manual dissection procedure was as follows. Flies were briefly anesthetized on ice and placed with fine forceps (Fine Science Tools) into a custom holder machined from black Delrin (Autotiv or Protolabs). This holder had an inverted pyramid shape to reduce obstruction of the eyes. The head was angled slightly forward to make the posterior side easier to access with the microscope objective. After clipping the wings, the head and thorax were fixed to the holder using UV-curable glue (Aquaseal UV Field Repair Adhesive), which was hardened with a short pulse of ultraviolet light (LED-200, Electro-Lite Co.). To limit large brain movements, the proboscis was also immobilized with a small amount of the same glue. A 30G\,$\times$\,1/2'' needle (BD PrecisionGlide Needle, 30G\,$\times$\,1/2'' Regular Wall Hypodermic, 305106) was then used in extracellular \textit{Drosophila} saline to cut an opening in the head capsule, and tracheoles and fat were removed to expose the brain. Brain motion was further minimized by stretching muscle 16 through gentle pulling on the oesophagus, or alternatively by removing the muscle via an anterior clip. The extracellular saline contained 103~mM NaCl, 3~mM KCl, 5~mM TES, 8~mM trehalose, 10~mM glucose, 26~mM NaHCO$_3$, 1~mM NaH$_2$PO$_4$, 1.5~mM CaCl$_2$, and 4~mM MgCl$_2$, with an osmolarity of 270--275~mOsm. The saline was oxygenated with carbogen (95\%~O$_2$, 5\%~CO$_2$), resulting in a final pH of approximately 7.3.

\subsection{Visual panorama and visual stimuli}

All calcium imaging and LED-based optogenetics experiments were performed using a custom-built LED display system adapted from the G4 display platform developed by the M.~Reiser lab. Visual stimuli were presented on a circular panoramic arena composed of modular 16\,$\times$\,16 pixel LED panels, with each pixel subtending approximately 1.875\textdegree{} of the fly's visual field. The arena spanned 60\textdegree{} in elevation and 330\textdegree{} in azimuth, with the 30\textdegree{} sector directly behind the animal removed to allow space for ball illumination and tracking hardware. The panels used blue LEDs with a peak emission at 470~nm, selected to minimize spectral overlap with GCaMP fluorescence. For calcium imaging experiments, five layers of Rosco R381 gel filters were placed in front of the arena to further reduce spectral interference, and an additional diffuser layer was added to suppress reflections (SXF-0600, Snow White Light Diffuser, Decorative Films). The visual stimulus consisted of a bright vertical bar with positive contrast, 8 pixels wide (15\textdegree{} azimuth). In most of the arena, the bar extended across the full two-panel height, except in the region spanning $-165$\textdegree{} to $+165$\textdegree{} behind the fly, where only a single display panel was present and the bar height was reduced by half.

In closed-loop experiments, the azimuthal position of the cue was controlled by the fly's yaw rotation measured from the spherical treadmill (see ``Spherical treadmill and locomotion measurement''). A yaw gain of 0.7 was used in all experiments, such that the displacement of the visual cue corresponded to 0.7 times the ball's yaw displacement. In calcium imaging experiments, the cue was pseudorandomly displaced by $\pm 90$\textdegree{} at intervals ranging from 60 to 80~s. After each perturbation, the cue remained stationary for 2~s before closed-loop coupling resumed. Cue position was recorded throughout the experiment using analog output signals from the display panels, which were acquired together with other experimental signals at 20~kHz via a National Instruments PCIe-6363 data acquisition card. During offline analysis, these analog signals were converted into cue position in pixel coordinates.

All two-photon optogenetics experiments were carried out using a custom visual stimulation system implemented in ViRMEn (Virtual Reality MATLAB Engine). Visual stimuli were delivered with a DLP projector (Texas Instruments DLP3010EVM-LC) operating at 120~Hz. To minimize contamination of imaging signals, the projector's red and green LEDs were disabled, and optical band-pass filters (Semrock 480/17-25 and 482/18-25) were placed at the projector output. The projected image was reflected by a 45\textdegree{} first-surface mirror onto a panoramic cone-shaped screen, which covered 60\textdegree{} in elevation and 330\textdegree{} in azimuth, again with the 30\textdegree{} region directly behind the fly omitted to accommodate treadmill illumination and tracking hardware. Stimuli were geometrically corrected to match the curvature of the screen.

In these experiments, the visual stimulus was a bright vertical bar spanning 15\textdegree{} in azimuth on a dark background. Except during heading perturbation trials, this bar served as a celestial cue and was coupled in closed loop exclusively to the fly's rotational movements. The cue position was updated using a gain of 0.7 relative to the fly's movement, consistent with prior work showing that this value improves heading tracking. During heading perturbations, the bar was randomly shifted by 90\textdegree{} to the left or right of its previous position and held at the new location for 2~s before closed-loop control was reinstated. All visual stimuli were controlled through ViRMEn.

\subsection{Experimental trial structure}

Before data collection, flies were acclimated for at least 15~min in closed loop with the visual cue. For virtual plume-tracking experiments, behavior was recorded in 16--18~min sessions. During each session, the fly remained in closed loop with the cue, which was displaced to a new angular position relative to its current location at random intervals of 60--80~s, alternating between $+90$\textdegree{} and $-90$\textdegree{}.

Two related stimulus designs were used, termed Protocol A and Protocol B. In both protocols, odor delivery was contingent on heading: odor was presented only when the fly's heading fell within a predefined angular sector relative to the visual landmark. Each session began with an initial blank period of 200~s. In Protocol A, this blank period was followed immediately by two odor blocks, yielding a three-block sequence (blank 1, odor zone 1, odor zone 2) with durations of 200, 400, and 400~s, respectively. In Protocol B, the session comprised a five-block sequence of alternating blank and odor epochs (blank 1, odor zone 1, blank 2, odor zone 2, blank 3). For Protocol B, two timing schedules were used: 200, 300, 150, 300, 150~s or 200, 350, 60, 350, 60~s, corresponding to the five successive blocks. Thus, in both protocols, odor was presented twice per session, first in one heading-defined sector and then in the opposite sector; Protocol A omitted an intervening blank period between the two odor blocks, whereas Protocol B included a brief clean-air interval between them. This design preserved a common experimental logic while varying block structure and duration to reduce stereotyped temporal expectation.

For calcium imaging experiments, trials were separated by 1~min of darkness. Because these experiments depended on spontaneously expressed behavior, trials were continued until the fly stopped walking.

\subsection{Definition of odor bout}

For analysis, the odor command signal was represented as a binary time series indicating valve OFF and ON states at each time point. Odor bouts were identified by first detecting individual ON epochs from OFF-to-ON and ON-to-OFF transitions in this signal. Because behaviorally contingent odor delivery often generated rapid switching of the valve during a single encounter, adjacent ON epochs separated by gaps of no more than 1~s were merged and treated as a single odor bout. Bout onset was defined as the start of the first ON epoch in the merged sequence, and bout offset as the end of the last ON epoch. For each bout, we measured both the overall bout duration, including any brief intervening OFF periods, and the cumulative odor-on duration, defined as the summed duration of all ON epochs within the bout. Bouts were included in subsequent analyses only if cumulative odor-on duration was at least 0.5~s.

\subsection{Computing exponentially weighted history metric}

To quantify recent odor experience at any time point during a trial, we computed an event-based exponentially weighted odor history variable. Odor encounter times were defined as transitions of the commanded valve state into the odor-on condition. For any time $t$, we identified all prior odor-on transition times $\{t_i\}$ satisfying $t_i < t$, and calculated odor history as $\sum_i \exp(-(t - t_i)/\tau)$, where $\tau$ denotes the exponential decay time constant. Under this formulation, recent odor encounters contribute more strongly than earlier encounters, with each contribution decaying continuously as a function of time since encounter onset. Because this measure is based on odor-on transition events rather than odor occupancy, it reflects the recent temporal structure and frequency of odor encounters, but does not scale with encounter duration. For bout-based analyses, this quantity was evaluated at the bout start time.

\subsection{Definition of sparse vs.\ dense regime}

To quantify locomotor behavior as a function of recent odor experience, we computed an event-based exponentially weighted odor-history variable continuously throughout each trial from the sequence of prior odor-on transitions using an exponential decay constant $\tau$. For each fly, the distribution of odor-history values across all trial time points was examined as a histogram, and clustered regions of this distribution were used to define odor-history regimes. The lowest-history regime, corresponding to periods before odor-history accumulation, was excluded from analysis. The highest-history regime was defined as the dense regime, reflecting periods with clustered recent odor encounters, and the intermediate regime was defined as the sparse regime, reflecting periods with intermittent but relatively infrequent recent odor encounters. All trial time points were labeled according to these regimes. This time-resolved segmentation was used to compute average locomotor statistics, including forward velocity, yaw speed, and trajectory straightness, during sparse and dense periods across the full trial. For bout-aligned analyses, each bout was assigned the label of the odor-history regime present at its start time, and these labels were used to compare locomotor response dynamics between sparse and dense odor-history conditions.

\subsection{Immunohistochemistry}

Brains from female flies 1--3 days posteclosion were dissected in \textit{Drosophila} external saline and fixed in 4\% paraformaldehyde (Electron Microscopy Sciences, 15714) in phosphate-buffered saline (PBS; Thermo Fisher Scientific, 46-013-CM) for 15~min at room temperature. Samples were washed in PBS and blocked for 20~min in PBS containing 5\% normal goat serum (Sigma-Aldrich, G9023) and 0.44\% Triton X-100 (Sigma-Aldrich, T8787). Brains were incubated in primary antibody diluted in blocking solution for approximately 24~h at room temperature, washed in PBS, and then incubated in secondary antibody in blocking solution for approximately 24~h at room temperature. Antibodies were selected according to the experimental protocol. After a final PBS wash, brains were mounted in antifade medium (Vectashield, Vector Laboratories, H-1000) for imaging. Images were acquired on a Leica SPE confocal microscope using a 40$\times$, 1.15~NA oil-immersion objective. Image stacks contained 50--200 z-sections collected at 1~\textmu m intervals and were recorded at a resolution of 1,024\,$\times$\,1,024 pixels. For visualization of Gal4 expression patterns, primary antibodies were chicken anti-GFP (1:1,000; Abcam, ab13970) and mouse anti-Bruchpilot (1:30; Developmental Studies Hybridoma Bank, nc82). Secondary antibodies were Alexa Fluor 488 goat anti-chicken (1:250; Invitrogen, A11039) and Alexa Fluor 633 goat anti-mouse (1:250; Invitrogen, A21050). For visualization of cell fills after whole-cell patch-clamp recordings, streptavidin-Alexa Fluor 568 (1:1,000; Invitrogen, S11226) was included in both the primary and secondary antibody solutions.

\subsection{Processing calcium imaging data}

Data analysis was carried out in either MATLAB 2020a or MATLAB R2023b. The calcium imaging dataset included X flies expressing GCaMP driven by MBON21-Gal4, X flies expressing GCaMP driven by MBON09-Gal4, X flies expressing GCaMP driven by the hDeltaB split-Gal4 line, X flies expressing GCaMP driven by the FB4R split-Gal4 line, X flies expressing GCaMP driven by the hDeltaB \& MBON09 split-Gal4 line, and X flies expressing GCaMP driven by the hDeltaB \& MBON21 split-Gal4 line. For each trial, rigid motion in the $x$, $y$, and $z$ dimensions was corrected using the NoRMCorre algorithm. Regions of interest (ROIs) were defined across the z-stack. For each ROI, $\Delta F/F$ was computed, with baseline fluorescence ($F$) taken as the mean of the lowest 10\% of fluorescence values within that trial, which lasted 600~s.

For imaging of the fan-shaped body, ROIs of approximately equal width were placed along the horizontal axis of the structure without overlapping. For MBON imaging, two ROIs were used: one corresponding to the left MBON and the other to the right MBON.

\subsection{Processing locomotion and visual arena data}

Spherical treadmill position was tracked online with machine-vision software (Fictrac v.2.1) and recorded as a voltage signal. For offline analysis, this voltage trace was converted to radians and unwrapped. The resulting signals were low-pass filtered with a second-order Butterworth filter with a 0.003 corner frequency, then downsampled to half of the ball-tracking update rate. Velocity was computed with the MATLAB \texttt{gradient} function. Spurious high velocity values exceeding 20~rad\,s$^{-1}$ were clipped to 20~rad\,s$^{-1}$, after which the timeseries were smoothed in MATLAB using the \texttt{smooth} function with the loess method and a 33~ms window, and then resampled to 60~Hz, matching the ball-tracking update rate. Forward and lateral velocities were converted to millimetres per second, while yaw velocity was expressed in degrees per second.

During calcium imaging, a signal from the imaging software marking the end of each volumetric stack was recorded on the same acquisition card as the online ball-tracking signals. The ball-kinematic data were then resampled to the imaging acquisition rate ($\sim$4--7~Hz depending on the ROI), enabling alignment with the acquired imaging volumes.

\subsection{Virtual plume paradigm simulation}

Fly trajectories were generated in discrete time with step size \(\Delta t\), inferred from the median sampling interval of the reference dataset. Reference forward velocity, lateral velocity, and yaw velocity were each modeled as first-order autoregressive processes. For a generic kinematic variable \(z_t\), the latent fluctuation dynamics were written as
\[
z_{t+1} = a_z z_t + \sigma_z \varepsilon_t,\qquad \varepsilon_t \sim \mathcal{N}(0,1),
\]
where \(a_z\) denotes the autoregressive coefficient estimated from the reference trajectory and \(\sigma_z\) denotes the innovation scale estimated from the corresponding residuals. Separate parameters \((a_{\mathrm{fw}},\sigma_{\mathrm{fw}})\), \((a_{\mathrm{side}},\sigma_{\mathrm{side}})\), and \((a_{\mathrm{yaw}},\sigma_{\mathrm{yaw}})\) were fit for forward, lateral, and yaw dynamics, respectively.

Heading dynamics were governed by an emergent-goal alignment rule defined from recent heading history. At simulation step \(t\), the model considered the set of headings observed within a lookback window of duration \(T_{\mathrm{win}}\), and computed the circular mean direction \(\mu_t\) together with the mean resultant length \(R_t \in [0,1]\),
\[
\mu_t = \mathrm{atan2}\!\left(\frac{1}{n_t}\sum_{k}\sin\theta_k,\;\frac{1}{n_t}\sum_{k}\cos\theta_k\right),
\qquad
R_t = \sqrt{\left(\frac{1}{n_t}\sum_{k}\cos\theta_k\right)^2+\left(\frac{1}{n_t}\sum_{k}\sin\theta_k\right)^2},
\]
where the sums run over the \(n_t\) samples in the recent window. After a minimum number of window samples, \(N_{\min}\), had been accumulated, this recent circular mean served as the instantaneous heading target. The corresponding alignment torque was
\[
\tau_t = K_t R_t \sin(\mu_t-\theta_{t-1}),
\]
where \(\theta_{t-1}\) is the previous heading and \(K_t\) is the time-dependent alignment gain. The factor \(R_t\) adaptively scaled the alignment term according to the directional coherence of recent headings, such that highly consistent recent motion generated stronger self-alignment than diffuse heading histories.

Yaw velocity was then updated as the sum of the autoregressive yaw fluctuation and the emergent-goal alignment torque, followed by odor-dependent gain modulation,
\[
\omega_t = \left(y_t + \tau_t\right) G^{\mathrm{yaw}}_t,
\]
where \(y_t\) denotes the AR(1) yaw state and \(G^{\mathrm{yaw}}_t\) is a multiplicative yaw-gain term determined by the current odor-response state. Heading was integrated according to
\[
\theta_t = \mathrm{wrap}_{2\pi}\!\left(\theta_{t-1} + \omega_t \Delta t\right).
\]

Forward velocity was modeled as a mean-reverting process around a prescribed mean speed \(V_{\mathrm{fw}}\),
\[
v^{\mathrm{fw}}_t = V_{\mathrm{fw}} + f_t + b^{\mathrm{fw}}_t,
\]
where \(f_t\) is the AR(1) forward-velocity state and \(b^{\mathrm{fw}}_t\) is a transient odor-dependent forward bias. Lateral velocity was generated as
\[
v^{\mathrm{side}}_t = s_t G^{\mathrm{side}}_t,
\]
where \(s_t\) is the AR(1) lateral-velocity state and \(G^{\mathrm{side}}_t\) is a transient odor-dependent lateral gain. Thus, both translational velocity components inherited temporal correlations from the reference kinematics while remaining modulated by odor encounter history.

Odor availability was defined by a temporal gate and an angular sector. After a specified onset time \(t_{\mathrm{odor}}\), odor was considered present whenever the instantaneous heading \(\theta_t\) lay within the angular interval \([\phi_{\mathrm{start}},\phi_{\mathrm{end}}]\). This produced a binary odor state \(O_t\), from which odor-entry events were identified as transitions from \(O_{t-1}=0\) to \(O_t=1\). Recent odor exposure was summarized by an exponentially decaying history variable \(H_t\),
\[
H_t = \lambda_H H_{t-1} + O_{t-1}\Delta t,
\qquad
\lambda_H = \exp\!\left(-\frac{\Delta t}{\tau_H}\right),
\]
where \(\tau_H\) is the odor-history time constant. Odor-response variables for yaw gain, lateral gain, and forward bias evolved as decaying internal states with recovery timescale \(\tau_R\). At each odor-entry event, the current odor history was compared with thresholds \(H_{\mathrm{sparse}}\) and \(H_{\mathrm{dense}}\) to classify the encounter regime, and regime-specific increments were applied to the yaw-gain, lateral-gain, and forward-bias states. These internal response states then relaxed exponentially over time and multiplicatively or additively modulated the kinematic updates described above. The alignment gain \(K_t\) was held constant throughout the simulation and was not modified after odor onset.

Finally, body-frame motion was transformed into world coordinates and integrated to obtain position. With heading \(\theta_t\), forward velocity \(v^{\mathrm{fw}}_t\), and lateral velocity \(v^{\mathrm{side}}_t\), position updates were computed as
\[
x_t = x_{t-1} + \left(v^{\mathrm{fw}}_t\cos\theta_t - v^{\mathrm{side}}_t\sin\theta_t\right)\Delta t,
\]
\[
y_t = y_{t-1} + \left(v^{\mathrm{fw}}_t\sin\theta_t + v^{\mathrm{side}}_t\cos\theta_t\right)\Delta t.
\]
This procedure produced simulated trajectories whose local kinematic fluctuations matched the temporal structure of the reference data, while global directional organization emerged from recent heading history and odor-dependent modulation.
